% !TEX root = chapter-standalone.tex

\section{Functorial stacking categories}
\label{sec:functorial-stacking}

\begin{ctdefinition}[Functorial stacking category]
    \label{def:functorial-stacking-cat}
    A \maindef{functorial stacking category} is a stacking \SY{category} where the two stacking operations~$\mtimescatob$ and~$\mtimescatmor$ are the two components of a \SY{functor} \begin{equation}
        \mtimescat\colon \CatC \Ctimes \CatC \to \CatC.
    \end{equation}
    In infix notation, this means that
    \begin{equation}
        (\mora \mthen \morc )
        \mtimescatmor (\morb \mthen \mord )
        =
        (\mora \mtimescatmor \morb )
        \mthen (\morc \mtimescatmor \mord )
    \end{equation}
    for all morphisms $\mora, \morb, \morc$ and $\mord$ (where respectively $\mora$ and $\morc$, and $\morb$ and $\mord$ are composable), and that
    \begin{equation}\label{eq:functorial-stacking-condition-identities}
        \catidat{\Obja} \mtimescatmor \catidat{\Objb} = \catidat{\Obja \mtimescatob \Objb}
    \end{equation}
    for all objects $\Obja$, $\Objb$ of $\CatC$.
\end{ctdefinition}

This describes a sort of commutativity property: we can either first compose horizontally and then vertically, or vice versa; either way, we obtain the same result~(\cref{fig:stacking-functorial}).

\begin{example}
    The cartesian product of sets and functions is a functorial stacking operation.
    Suppose we are given functions $\mora, \morb, \morc$ and $\mord$ (where respectively $\mora$ and $\morc$, and $\morb$ and $\mord$ are composable).
    On the one hand,
    \begin{align}
        ((\mapa \mthen \mapb) \funcprod (\mapc \mthen \mapd))(\tup{\setAel, \setBel}) & = \tup{(\mapa \mthen \mapb)(\setAel), (\mapc \mthen \mapd)(\setBel)} \\
                                                                                      & = \tup{\mapb(\mapa(\setAel)), \mapd(\mapc(\setBel))}
    \end{align}
    while on the other hand,
    \begin{align}
        ((\mapa \funcprod \mapc) \mthen (\mapb \funcprod \mapd))(\tup{\setAel, \setBel}) & = (\mapb \funcprod \mapd)((\mapa \funcprod \mapc)(\tup{\setAel, \setBel})) \\
                                                                                         & = (\mapb \funcprod \mapd)(\tup{\mapa(\setAel),  \mapc(\setBel)}) \\
                                                                                         & = \tup{\mapb(\mapa(\setAel)), \mapd(\mapc(\setBel))}
    \end{align}
\end{example}

\begin{example}
    The sum of sets and functions is a functorial stacking operation.
\end{example}

\todotext{J: product and sum for Rel -- make exercise or graded exercise perhaps?}

\begin{example}
    Consider the stacking operation defined previously for the category of real matrices.
    It is functorial stacking:
    \begin{align}
        (\mat{A} \mthen \mat{B}) \mtimescatmor (\mat{C} \mthen \mat{D}) & = (\mat{B}\mat{A}) \mtimescatmor (\mat{D}\mat{C}) \\
                                                                        & =
        \begin{bmatrix}
            \mat{BA} & \mat{0}  \\
            \mat{0}  & \mat{DC}
        \end{bmatrix}
    \end{align}
    \begin{align}
        (\mat{A} \mtimescatmor \mat{C}) \mthen (\mat{B} \mtimescatmor \mat{D}) & =
        \begin{bmatrix}
            \mat{A} & \mat{0} \\
            \mat{0} & \mat{C}
        \end{bmatrix}
        \mthen
        \begin{bmatrix}
            \mat{B} & \mat{0} \\
            \mat{0} & \mat{D}
        \end{bmatrix} \\
                                                                               & =
        \begin{bmatrix}
            \mat{B} & \mat{0} \\
            \mat{0} & \mat{D}
        \end{bmatrix}
        \begin{bmatrix}
            \mat{A} & \mat{0} \\
            \mat{0} & \mat{C}
        \end{bmatrix} \\
                                                                               & =
        \begin{bmatrix}
            \mat{BA} & \mat{0}  \\
            \mat{0}  & \mat{DC}
        \end{bmatrix}
    \end{align}
\end{example}

\begin{example}
    The stacking operation previously defined for the category $\DP$ of design problems is functorial.
\end{example}

\subsection{Stacking for $\LTI$ is almost functorial}

\begin{example}
    We want to show that \LTI, equipped with the defined stacking operations, is almost a \SY{functorial stacking semicategory}, but not quite.
    Given morphisms~$\mora\colon \styleobj{l}\mto \styleobj{m}$,~$\morc\colon \styleobj{m}\mto \styleobj{n}$,~$\morb\colon \styleobj{o}\mto \styleobj{p}$,~$\mord\colon \styleobj{p}\mto \styleobj{q}$, we would need to have
    \begin{equation*}
        (\mora \mthen \morc)
        \mtimescatmor (\morb\mthen \mord)=(\mora\mtimescatmor \morb)\mthen (\morc\mtimescatmor \mord).
    \end{equation*}
    This, however, is not true.
    Let's see this by looking at the first matrix component of the LTI system.
    On one hand we have:
    \begin{equation}\label{eq:checking-LTI-functorial-stacking-1}
        \mat{A}_{(\mora \mthen \morc)\mtimescatmor (\morb\mthen \mord)}=
        \begin{bmatrix}
            \mat{A}_\mora              & \mat{0}       & \mat{0}                    & \mat{0}       \\
            \mat{B}_\morc\mat{C}_\mora & \mat{A}_\morc & \mat{0}                    & \mat{0}       \\
            \mat{0}                    & \mat{0}       & \mat{A}_\morb              & \mat{0}       \\
            \mat{0}                    & \mat{0}       & \mat{B}_\mord\mat{C}_\morb & \mat{A}_\mord
        \end{bmatrix}.
    \end{equation}
    On the other hand we have:
    \begin{equation}\label{eq:checking-LTI-functorial-stacking-2}
        \mat{A}_{(\mora\mtimescatmor \morb)\mthen (\morc\mtimescatmor \mord)}=
        \begin{bmatrix}
            \mat{A}_\mora               & \mat{0}                     & \mat{0}       & \mat{0}       \\
            \mat{0}                     & \mat{A}_\morb               & \mat{0}       & \mat{0}       \\
            \mat{B}_\morc \mat{C}_\mora & \mat{0}                     & \mat{A}_\morc & \mat{0}       \\
            \mat{0}                     & \mat{B}_\mord \mat{C}_\morb & \mat{0}       & \mat{A}_\mord
        \end{bmatrix}.
    \end{equation}
    These two are different, and will therefore describe different systems.
    However, the two matrices just differ by two permutations, which can be expressed via an invertible linear transformation~$\mat{T}$ as follows:
    \begin{equation}\label{eq:checking-LTI-functorial-stacking-3}
        \mat{A}_{(\mora\mtimescatmor \morb)\mthen (\morc\mtimescatmor \mord)}=
        \underbrace{\begin{bmatrix}
                1 & 0 & 0 & 0 \\
                0 & 0 & 1 & 0 \\
                0 & 1 & 0 & 0 \\
                0 & 0 & 0 & 1
            \end{bmatrix}}_{\mat{T}}\cdot
        \mat{A}_{(\mora \mthen \morc)\mtimescatmor (\morb\mthen \mord)}
        \cdot
        \begin{bmatrix}
            1 & 0 & 0 & 0 \\
            0 & 0 & 1 & 0 \\
            0 & 1 & 0 & 0 \\
            0 & 0 & 0 & 1
        \end{bmatrix}.
    \end{equation}
    It turns out that $(\mora \mthen \morc)\mtimescatmor (\morb\mthen \mord)$ and~$(\mora\mtimescatmor \morb)\mthen (\morc\mtimescatmor \mord)$ are equivalent systems (\cref{def:equivalence_lti}), even though they are not equal.
    In particular, although \LTI is not a \SY{functorial stacking semicategory}, by \cref{lem:equivsystequivact}~$(\mora \mthen \morc)\mtimescatmor (\morb\mthen \mord)$ and~$(\mora\mtimescatmor \morb)\mthen (\morc\mtimescatmor \mord)$ have the same action, and hence \LTI is, in one sense, ``morally'' functorial.
\end{example}

\begin{gradedexercise}[\exname{StackingLTI}]
    \label{ex:StackingLTI}

    Consider the category of finite-dimensional linear time-invariant systems defined in \linkifexists{def:LTICat} with the stacking defined above.

    Your task:
    supposing that morphisms~$\mora\colon \styleobj{l}\mto \styleobj{m}$,~$\morc\colon \styleobj{m}\mto \styleobj{n}$,~$\morb\colon \styleobj{o}\mto \styleobj{p}$,~$\mord\colon \styleobj{p}\mto \styleobj{q}$ are given, compute the matrices
    \begin{equation}
        \mat{A}_{(\mora \mthen \morc)\mtimescatmor (\morb\mthen \mord)}
    \end{equation}
    and
    \begin{equation}
        \mat{A}_{(\mora\mtimescatmor \morb)\mthen (\morc\mtimescatmor \mord)}
    \end{equation}
    associated with the morphisms $(\mora \mthen \morc)\mtimescatmor (\morb\mthen \mord)$ and $(\mora\mtimescatmor \morb)\mthen (\morc\mtimescatmor \mord)$, respectively.
\end{gradedexercise}

\solutionof{StackingLTI}

\begin{marginfigure}
    \centering
    $(\mora \mtimescatmor \morb )
        \mthen (\morc \mtimescatmor \mord )$

    \vspace{3mm}
    \includesag{commuting_stacking_eff}

    $=$

    \vspace{3mm}
    \includesag{commuting_stacking_eff_bis}
    %\includegraphics[width=6.5cm]{effects8}
    \vspace{3mm}
    $(\mora \mthen \morc ) \mtimescatmor (\morb \mthen \mord )$

    \caption{Commutation of stacking and composition in a functorial stacking \SY{semicategory}. }
    \label{fig:stacking-functorial}
\end{marginfigure}

%More in detail, we have that for any~$n \geq 2$, there is a semifunctor
%%
%\begin{equation}
%    \mtimescat^{(n)}\colon \underbrace {\CatC \times \CatC \times \cdots \times \CatC}_{\text{$n$ times}} \to \CatC.
%\end{equation}
%

\subsection{A very non-functorial stacking}

We describe here a \SY{category} that will serve as an example of a stacking \SY{category} that is not functorial.

There are two types of computation: on the one hand, ``functional'', or ``pure'', or ``free of side effects'', and on the other hand there is \emph{effectful} computation.

In functional programming, functions in the programming language are ``pure'' and are very much like mathematical functions: they need an input and produce and output.
And they don't interfere with other functions.

Effectful procedures, instead, can ``change the world'': for example, printing a page, sending an email, or placing an order of pizza.
The order in which effectful procedures are evaluated might change the result.
For example, the result of the sequence of operations
%
\begin{enumerate}
    \item Order a pizza;
    \item Cancel the last order;
\end{enumerate}
%
is different from the result of the sequence of operations
%
\begin{enumerate}
    \item Cancel the last order;
    \item Order a pizza;
\end{enumerate}

A very elegant way to treat side effects mathematically is using \emph{linear types}~\cite{Wadler90lineartypes}.
We will mention those in a successive part on linear logic.
For now, we stick to a simple treatment.

\begin{marginfigure}
    \centering
    \subfloat[\label{fig:effects1}
        A morphism~$\mora\colon \Obja \mtoin{\Effects} \Objb$ in \Effects.
    ]{
        \rule{1cm}{0pt}\includesag{morphism_eff}\rule{1cm}{0pt}
    }\\
    \subfloat[\label{fig:effects2}
        Its representation in \SetL as a morphism
        $\repreff{\mora}\colon \cObj{\Obja,\EfW} \mtoin{\SetL} \cObj{\Objb, \EfW}$.
    ]{
        \rule{1cm}{0pt}\includesag{morphism_tup}\rule{1cm}{0pt}
    }
    \caption{}
    \label{fig:effects12}
\end{marginfigure}

We are going to define a category \Effects.
The idea is to add another variable that represents ``the world'' that can be affected.
An effectful function
%
\begin{equation}\label{eq:Eff-morphism}
    \mora\colon \Obja \mtoin{\Effects} \Objb,
\end{equation}
%
which could have some unknown side effects on the world, can be represented by a pure function
\begin{equation}
    \repreff{\mora}\colon \Obja \cartprod \EfW \mtoin{\Set} \Objb \cartprod \EfW,
\end{equation}\label{eq:Eff-rep-morphism-Set}
or, in other words, as a morphism
\begin{equation}\label{eq:Eff-rep-morphism-SetL}
    \repreff{\mora}\colon \cObj{\Obja,\EfW} \mtoin{\SetL} \cObj{\Objb, \EfW},
\end{equation}
%
where~$\EfW$ is the set of all possible worlds~(\cref{fig:effects12}).

The second input to~$\repreff{\mora}$ is the state of the world before the execution of the function.
The second output of~$\repreff{\mora}$ is the state of the world after the execution of the function.

We can now interconnect different effectful functions, with some precautions.
We cannot ``split the world'', by creating a function of type~$\EfW \to \EfW \cartprod \EfW$.
We will re-state this formally when we get to linear logic.

\begin{marginfigure}
    \centering
    \subfloat[\label{fig:effects3}
        Composition in \Effects.
    ]{
        \includesag{morphism_eff_comp}
        %\includegraphics[width=4.5cm]{effects3}
    }\\
    \subfloat[\label{fig:effects3b}
        Its representation in \SetL.
    ]{
        \includesag{morphism_tup_comp}
        %\includegraphics[width=4.5cm]{effects3b}
    }
    \caption{}
    \label{fig:effect-composition}
\end{marginfigure}

We can extend usual function composition to composition of effectful functions as in~\cref{fig:effect-composition}.
The second effectful function operates on the world after it was possibly modified by the first effectful function.

We have all the ingredients to define the category \Effects of effectful computation.

\begin{definition}[Category of effectful procedures \Effects]
    \SYNDEF{semicategory of sets and functions with side effects}
    \label{def:Effects}
    Fix a set~$\EfW$ of all possible worlds.
    The category~$\Effects$ is defined by the following:
    \begin{itemize}
        \item \emph{Objects}: same as the objects of \SetL;
        \item \emph{Morphisms}: a morphism~$\mora\colon \Obja \mtoin{\Effects} \Objb$ is a function
              \begin{equation}
                  \repreff{\mora}\colon \cObj{\Obja, \EfW} \mtoin{\SetL} \cObj{\Objb, \EfW}.
              \end{equation}
        \item \emph{Composition}: The composition of~$\mora\colon \Obja \mtoin{\Effects} \Objb$ and~$\morb\colon \Objb \mtoin{\Effects} \Objc$ is the morphism~$\morab\colon \Obja \mtoin{\Effects} \Objc$ with~$\repreff{\morab}$ given by
              \begin{equation}
                  \repreff{\mora \mthenof \Effects \morb}
                  = \repreff{\mora} \mthenof{\SetL} \repreff{\morb}
              \end{equation}
              as illustrated in \cref{fig:effect-composition}.
        \item \emph{Identities}: for any object $\Obja$, its identity morphism $\idmon_\Obja \colon \Obja \mtoin{\Effects} \Obja$ is the identity function
              \begin{equation}
                  \cObj{\Obja, \EfW} \mtoin{\SetL} \cObj{\Obja, \EfW}.
              \end{equation}
    \end{itemize}
\end{definition}

\begin{marginfigure}
    \centering
    \subfloat[\label{fig:effects4}
        Stacking in \Effects]{
        \includesag{stacking_eff}
        %\includegraphics[width=4.5cm]{effects4}
    }

    \subfloat[\label{fig:effects5} Its representation in \SetL]{
        %\includegraphics[width=4.5cm]{effects5}
        \includesag{stacking_tup}
    }
    \caption{}
    \label{fig:effect-stacking}
\end{marginfigure}

We can now make \Effects into a stacking \SY{category} by deciding how to evaluate a stack of functions.
We define this in the way shown in~\cref{fig:effect-stacking}.

\begin{lemma}
    \label{lem:effects-not-functorial}
    \Effects is \emph{not} a \SY{functorial stacking semicategory}.
\end{lemma}

\begin{proof}
    In general, we have
    \begin{equation}
        (\mora \mtimescatmor \morb )
        \mthen (\morc \mtimescatmor \mord )
        \neq
        (\mora \mthen \morc ) \mtimescatmor (\morb \mthen \mord ).
    \end{equation}
    This is shown graphically in \cref{fig:effects-non-functorial}.
\end{proof}

\begin{figure*}[h]
    \centering
    \hfill
    \subfloat[\label{fig:effects7}$(\mora \mtimescatmor \morb )
            \mthen (\morc \mtimescatmor \mord )$]{
        %\includegraphics[width=6.5cm]{effects7}
        \includesag{stacking_tup_counterproof}
    }
    \hfill
    \subfloat[\label{fig:effects9} $(\mora \mthen \morc ) \mtimescatmor (\morb \mthen \mord )$]{
        %\includegraphics[width=6.5cm]{effects9}
        \includesag{stacking_tup_counterproof_bis}
    }
    \hfill
    \caption{Proof that \Effects is not a \SY{functorial stacking} category by showing that
        the two morphisms above have different representations in~\SetL.}
    \label{fig:effects-non-functorial}
\end{figure*}

\todotextjira{751}{\alphubel  @GZ: There is a typo in the above diagrams -- U and Z are switched from where they should be placed.
    Also, the right hand diagram should have *two* boxes with dotted lines, not just one.
}

% \begin{comment}
% \devel{
%
%     \begin{equation}
%         \begin{aligned}
%             \mora\colon \cObj{\setAn{1}, \ldots,\setAn{m}} & \mto \cObj{\setBn{1}, \ldots, \setBn{n}}, \\
%             \morb\colon \cObj{\setDn{1}, \ldots,\setDn{o}} & \mto \cObj{\setEn{1}, \ldots, \setEn{p}}, \\
%             \morc\colon \cObj{\setBn{1}, \ldots,\setBn{n}} & \mto \cObj{\setCn{1}, \ldots, \setCn{q}}, \\
%             \mord\colon \cObj{\setEn{1}, \ldots,\setEn{p}} & \mto \cObj{\setFn{1}, \ldots, \\setFn{r}},
%         \end{aligned}
%     \end{equation}
%     to which maps~$\mape,\mapf,\mapg,\maph$ are associated.
%
%     We want to show that
%     \begin{equation}
%         (\mora \mtimescat \morb )
%         \mthen (\morc \mtimescat \mord )
%         =
%         (\mora \mthen \morc ) \mtimescat (\morb \mthen \mord ),
%     \end{equation}
%     and we do it by showing that the map associated with the left-hand side of the equation corresponds to the one on the right-hand side.
%     We know the maps associated to~$(\mora \mtimescat \morb )$ and~$(\morc \mtimescat \mord )$, which are
%     \begin{equation}
%         \cohm_{\setA,\setD}\then (\mape \funcprod \mapg)\then \cohm_{\setB,\setE},
%     \end{equation}
%     and
%     \begin{equation}
%         \cohm_{\setB,\setE}\then (\mapf \funcprod \maph)\then \cohm_{\setC,\setF},
%     \end{equation}
%     respectively.
%     The map associated to~$(\mora \mtimescat \morb )
%         \mthen (\morc \mtimescat \mord )$ is therefore
%     \begin{equation}
%         \label{eq:setL_functorial_a}
%         \cohm_{\setA,\setD}\then (\mape \funcprod \mapg)\then \cohm_{\setB,\setE}\then \cohm_{\setB,\setE}\then (\mapf \funcprod \maph)\then \cohm_{\setC,\setF}.
%     \end{equation}
%     On the other hand, the map associated to~$(\mora \mthen \morc )$ is~$\mape\then \mapf$, and the map associated to~$(\morb \mthen \mord)$ is~$\mapg\then \maph$.
%     Therefore, the map associated to~$(\mora \mthen \morc ) \mtimescat (\morb \mthen \mord )$ is
%     \begin{equation}
%         \cohm_{\setA,\setD}\then ((\mape\then \mapf)\funcprod (\mapg\then \maph))\then \cohm_{\setC,\setD}
%     \end{equation}
%
%     \todotext{Now need to finish by showing equivalence of maps}
% }
% \end{comment}
